% Titolo: Primo incontro conoscitivo
\documentclass[a4paper, 11pt]{article}

% --- PACCHETTI ---
\usepackage[utf8]{inputenc}
\usepackage[italian]{babel}
\usepackage[margin=2.5cm]{geometry}
\usepackage{graphicx}
\usepackage{tabularx}
\usepackage{hyperref}
\usepackage{fancyhdr}
\usepackage{array}
\usepackage{titlesec}
\usepackage{xcolor}

% --- CONFIGURAZIONE HEADER E FOOTER ---
\pagestyle{fancy}
\fancyhf{}
\fancyhead[L]{Bit Crusher}
\fancyhead[R]{Verbale Interno N. 1} 
\fancyfoot[C]{\thepage}
\renewcommand{\headrulewidth}{0.4pt}

% Configurazione link indice
\hypersetup{
    colorlinks=true,
    linkcolor=black,
    urlcolor=blue
}

\begin{document}

% ==========================================
% FRONTESPIZIO
% ==========================================
\begin{titlepage}
    \centering
    
    % --- LOGHI ---
    % Sostituisci i \framebox con i comandi \includegraphics corretti
    \begin{minipage}{0.45\textwidth}
        \centering
        % Esempio: \includegraphics[width=0.7\textwidth]{logo_unipd.png}
        \includegraphics[width=0.7\textwidth]{unipd_logo.png}
    \end{minipage}
    \hfill
    \begin{minipage}{0.45\textwidth}
        \centering
        % Esempio: \includegraphics[width=0.7\textwidth]{WhatsApp Image 2026-10-15 at 17.29.12.jpeg}
        \includegraphics[width=0.7\textwidth]{Logo-BitCrusher-Nero.png}
    \end{minipage}

    \vspace{1.5cm}
    {\Large \textbf{UNIVERSITÀ DEGLI STUDI DI PADOVA}} \\
    \vspace{0.3cm}
    {\large Corso di Laurea in Informatica} \\
    \vspace{0.3cm}
    {\large Insegnamento di Ingegneria del Software} \\
    \vspace{0.3cm}
    {\large Anno Accademico 2026/2027} \\

    \vspace{2cm}
    {\Huge \textbf{Verbale di Riunione}} \\
    \vspace{0.5cm}
    {\Large Riunione Interna \#1}

    \vspace{2cm}
    \begin{flushleft}
        \large
        \textbf{Gruppo:} Bit Crusher \\
        \textbf{Email:} bitcrusherteam@gmail.com \\
        \textbf{Versione:} v1.0.0 \\
        \textbf{Data della riunione:} 2026-10-17 \\
        \textbf{Data ultima modifica:} \today
    \end{flushleft}
    \vfill
\end{titlepage}

\newpage

% ==========================================
% REGISTRO DELLE MODIFICHE
% ==========================================
\section*{Registro delle Modifiche}
\vspace{0.5cm}
\begin{center}
\renewcommand{\arraystretch}{1.5}
\begin{tabularx}{\textwidth}{|l|l|l|l|X|}
\hline
\textbf{Versione} & \textbf{Data} & \textbf{Autore} & \textbf{Verificatore} & \textbf{Descrizione} \\ \hline
1.0.0 & 2026-11-03 & Alex Singh & - & Approvazione per il rilascio. \\ \hline
0.2.2 & 2026-11-03 & Antonio Tufano & Riccardo Pozzato & Aggiornamento logo UniPD e adeguamento registro modifiche. \\ \hline
0.2.1 & 2026-10-19 & Marco Bollettin & Riccardo Pozzato & Revisione e modifiche della versione precedente \\ \hline
0.2.0 & 2026-10-17 & Antonio Tufano & Marco Bollettin & Aggiustamenti grafici del layout. \newline Stesura sezioni: OdG, resoconto e to-do. \\ \hline
0.1.0 & 2026-10-17 & Riccardo Pozzato & Antonio Tufano & Approvazione e prima stesura del verbale. \\ \hline


\end{tabularx}
\end{center}

\newpage

% ==========================================
% INDICE
% ==========================================
\tableofcontents
\newpage

% ==========================================
% 1. INFORMAZIONI GENERALI
% ==========================================
\section{Informazioni Generali}
\begin{itemize}
    \item \textbf{Luogo/Piattaforma:} Riunione Meet
    \item \textbf{Data:} 2026-10-17
    \item \textbf{Orario di inizio:} 17:30
    \item \textbf{Orario di fine:} 18:45
    \item \textbf{Responsabile:} Non ancora definito al momento della riunione
    \item \textbf{Scriba:} Riccardo Pozzato
\end{itemize}

\vspace{0.5cm}
\textbf{Partecipanti presenti:} \\
\vspace{0.2cm}
\renewcommand{\arraystretch}{1.2}
\begin{tabularx}{0.8\textwidth}{|X|c|}
\hline
\textbf{Nome e Cognome} & \textbf{Matricola} \\ \hline
Pietro Leonardo Acampora & 2146762 \\ \hline
Marco Bollettin & 2147955 \\ \hline
Alexandru Mitu & 2101083 \\ \hline
Riccardo Pozzato & 2147988 \\ \hline
Alex Singh & 2148621 \\ \hline
Antonio Tufano & 2137972 \\ \hline
\end{tabularx}

\vspace{0.5cm}
\textbf{Partecipanti assenti:} \\
\vspace{0.2cm}
\renewcommand{\arraystretch}{1.2}
\begin{tabularx}{0.8\textwidth}{|X|c|}
\hline
\textbf{Nome e Cognome} & \textbf{Matricola} \\ \hline
Erik Giuliano & 2138012 \\ \hline
\end{tabularx}

% ==========================================
% 2. ORDINE DEL GIORNO
% ==========================================
\vspace{0.5cm}
\section{Ordine del Giorno}
Gli argomenti previsti per la riunione odierna sono:
\begin{enumerate}
    \item Scelta del nome del gruppo.
    \item Scelta e approvazione del logo del gruppo.
    \item Creazione della mail ufficiale del gruppo.
    \item Creazione  della repository GitHub del gruppo e organizzazione delle directory.
    \item Discussione preliminare e scelta dei capitolati.
    \item Creazione dei canali di comunicazione interni ed esterni (WhatsApp, Discord).
    \item Scelta della tecnologia per la documentazione (LaTeX) e layout di impaginamento.
\end{enumerate}

% ==========================================
% 3. RESOCONTO DELLA RIUNIONE
% ==========================================
\newpage
\section{Resoconto della Riunione}

\subsection{Scelta del Nome del Gruppo}
Il gruppo ha scelto il nome \textbf{Bit Crusher}. Motivazione: identificare il team in modo univoco attraverso un brand affine all'ambito informatico.

\subsection{Logo del Gruppo}
È stato approvato all'unanimità il logo ufficiale del team, raffigurante un teschio in pixel art con un effetto digitale/glitch verde e azzurro. Verrà impiegato in tutti i documenti formali e sulle piattaforme condivise.

\subsection{Mail del Gruppo}
Per tutte le comunicazioni formali verso l'esterno, i docenti e le aziende proponenti, è stato creato il seguente indirizzo email ufficiale: \texttt{bitcrusherteam@gmail.com}.

\subsection{Git del Gruppo e Suddivisione Directory}
È stata creata l'organizzazione su GitHub e configurata la repository dedicata. La documentazione e il codice saranno suddivisi in specifiche directory (\texttt{docs/}, \texttt{src/}, ecc.) per mantenere il repository ordinato e facilitare il versionamento.

\subsection{Discussione sui Capitolati}
I membri si sono confrontati sui capitolati proposti valutando le tecnologie richieste. È stato definito un ordine di preferenza preliminare e si è deciso di contattare a breve le aziende per organizzare incontri conoscitivi di approfondimento. L'ordine di preferenza dei capitolati, concordato tramite un sondaggio interno, è il seguente:
\begin{itemize}
    \item \textbf{Prima Scelta:} Capitolato C6, presentato da VarGroup;
    \item \textbf{Seconda Scelta:} Capitolato C3, presentato da Uranion;
    \item \textbf{Terza Scelta:} Capitolato C1, presentato da Vimar.
\end{itemize}

\subsection{Canali di Comunicazione}
Sono stati ufficialmente adottati i seguenti strumenti di comunicazione:
\begin{itemize}
    \item \textbf{Gruppo WhatsApp:} per comunicazioni quotidiane e aggiornamenti rapidi.
    \item \textbf{Server Discord:} come piattaforma principale per le riunioni di gruppo, condivisione schermo e organizzazione in canali tematici.
    \item \textbf{Riunioni Meet/Zoom:} piattaforme per le riunioni esterne con aziende, per delucidazioni più articolate e allineamenti su come procede il lavoro.
    \item \textbf{Indirizzo Email:} per comunicazione veloci e concise con professori o aziende.
\end{itemize}

\subsection{Creazione Documentazione e Layout}
È stato deciso di adottare \textbf{LaTeX} per la stesura dell'intera documentazione di progetto, garantendo coerenza tipografica. È stato approvato il presente layout documentale per i verbali futuri.

\subsection{Decisioni prese}
\begin{itemize}
    \item \textbf{V-1.1} Il nome del gruppo sarà "Bit Crusher".
    \item \textbf{V-1.2} Approvazione del logo in stile pixel art glitch.
    \item \textbf{V-1.3} L'email ufficiale è stata definita in \texttt{bitcrusherteam@gmail.com}.
    \item \textbf{V-1.4} Istituzione dell'organizzazione GitHub e della repository documentale.
    \item \textbf{V-1.5} I canali di comunicazione ufficiali saranno WhatsApp, Discord, Meet/Zoom e email.
    \item \textbf{V-1.6} La documentazione verrà impaginata e templatizzata tramite LaTeX.
    \item \textbf{V-1.7} Concordati i principali capitolati sui quali il gruppo si concentrerà (C6, C3, C1)
\end{itemize}

% ==========================================
% 4. AZIONI DA SVOLGERE (TO-DO)
% ==========================================
\vspace{0.5cm}
\section{Azioni da Svolgere (To-Do)}
\begin{tabularx}{\textwidth}{|c|X|l|l|}
\hline
\textbf{Codice Rif.} & \textbf{Descrizione Task} & \textbf{Assegnatario} & \textbf{Scadenza} \\ \hline
VI-1.1 & Leggere documentazione riguardante i capitolati in autonomia e abbozzare possibili domande e chiarimenti da porre alle aziende. & Intero Gruppo & 2026-10-19 \\ \hline
VI-1.2 & Riorganizzare la repository GitHub per avere una prima versione organizzata & Marco Bollettin & 2026-10-19 \\ \hline
VI-1.3 & Pianificare seconda riunione interna tramite i canali di comunicazione concordati & Intero Gruppo & 2026-10-18 \\ \hline

\end{tabularx}

% ==========================================
% 5. PROSSIMA RIUNIONE
% ==========================================
\vspace{0.5cm}
\section{Prossima Riunione}
\begin{itemize}
    \item \textbf{Data:} 2026-10-21
    \item \textbf{Ora:} 17:30
    \item \textbf{OdG preliminare:} Stesura delle domande riguardanti i capitolati e primo contatto con le aziende tramite email, definizione ruoli del gruppo.
\end{itemize}

\end{document}